
\documentclass[12pt,a4paper]{article}

% ── Packages ─────────────────────────────────────────────────────────────
\usepackage[T1]{fontenc}
\usepackage[utf8]{inputenc}
\usepackage[margin=2.5cm]{geometry}
\usepackage{xcolor}
\usepackage{graphicx}
\usepackage{listings}
\usepackage{enumitem}
\usepackage{hyperref}
\usepackage{fancyhdr}
\usepackage{tcolorbox}
\usepackage{booktabs}
\usepackage{array}
\usepackage{tabularx}
\usepackage{titlesec}
\usepackage{mdframed}
\usepackage{fontawesome5}
\usepackage{lmodern}
\usepackage{microtype}
\usepackage{parskip}

% ── Colour Palette ─────────────────────────────────────────────────────
\definecolor{eyantraBlue}{HTML}{1E3A5F}
\definecolor{eyantraAccent}{HTML}{2D7DD2}
\definecolor{eyantaGreen}{HTML}{06A77D}
\definecolor{eyantraWarn}{HTML}{F4A261}
\definecolor{eyantraDanger}{HTML}{E63946}
\definecolor{codebg}{HTML}{F4F6F9}
\definecolor{codeborder}{HTML}{CDD4E0}
\definecolor{lightgray}{HTML}{F8F9FA}
\definecolor{darkgray}{HTML}{495057}

% ── Hyperlinks ─────────────────────────────────────────────────────────
\hypersetup{
  colorlinks=true,
  linkcolor=eyantraAccent,
  urlcolor=eyantraAccent,
  pdftitle={EEP Auto-Grading Platform -- Student Guide},
  pdfauthor={E-Yantra, IIT Bombay},
}

% ── Code Listing Style ─────────────────────────────────────────────────
\lstdefinestyle{shellstyle}{
  backgroundcolor=\color{codebg},
  basicstyle=\ttfamily\small,
  breaklines=true,
  frame=single,
  rulecolor=\color{codeborder},
  framesep=4pt,
  xleftmargin=6pt,
  xrightmargin=6pt,
  showstringspaces=false,
  keywordstyle=\color{eyantraBlue}\bfseries,
  commentstyle=\color{darkgray}\itshape,
  stringstyle=\color{eyantaGreen},
}
\lstset{style=shellstyle}

% ── Custom tcolorbox environments ──────────────────────────────────────
\tcbuselibrary{skins,breakable,listingsutf8}

\newtcolorbox{infobox}[1][Note]{
  colback=eyantraAccent!8,
  colframe=eyantraAccent,
  fonttitle=\bfseries,
  title={\faInfoCircle\enspace #1},
  breakable,
}

\newtcolorbox{warnbox}[1][Warning]{
  colback=eyantraWarn!12,
  colframe=eyantraWarn!80!black,
  fonttitle=\bfseries,
  title={\faExclamationTriangle\enspace #1},
  breakable,
}

\newtcolorbox{successbox}[1][Success]{
  colback=eyantaGreen!10,
  colframe=eyantaGreen,
  fonttitle=\bfseries,
  title={\faCheckCircle\enspace #1},
  breakable,
}

\newtcolorbox{dangerbox}[1][Important]{
  colback=eyantraDanger!8,
  colframe=eyantraDanger,
  fonttitle=\bfseries,
  title={\faBan\enspace #1},
  breakable,
}

% ── Section Title Styling ──────────────────────────────────────────────
\titleformat{\section}
  {\Large\bfseries\color{eyantraBlue}}
  {\thesection}{1em}{}[\color{eyantraAccent}\rule{\linewidth}{1.5pt}]

\titleformat{\subsection}
  {\large\bfseries\color{eyantraBlue}}
  {\thesubsection}{0.8em}{}

\titleformat{\subsubsection}
  {\normalsize\bfseries\color{darkgray}}
  {\thesubsubsection}{0.7em}{}

% ── Header / Footer ───────────────────────────────────────────────────
\pagestyle{fancy}
\fancyhf{}
\fancyhead[L]{\color{eyantraBlue}\small\bfseries EEP Auto-Grading Platform}
\fancyhead[R]{\color{darkgray}\small Student Guide}
\fancyfoot[C]{\color{darkgray}\small\thepage}
\renewcommand{\headrulewidth}{0.5pt}
\renewcommand{\headrule}{\color{eyantraAccent}\hrule width\headwidth height\headrulewidth}

% ── Step counter ──────────────────────────────────────────────────────
\newcounter{stepnum}
\newcommand{\step}[1]{%
  \stepcounter{stepnum}%
  \vspace{6pt}%
  \noindent\colorbox{eyantraBlue}{\color{white}\bfseries\small\ Step \thestepnum\ }%
  \hspace{4pt}\textbf{#1}\\[4pt]%
}

% ─────────────────────────────────────────────────────────────────────
\begin{document}
% ─────────────────────────────────────────────────────────────────────

% ── Title Page ────────────────────────────────────────────────────────
\begin{titlepage}
  \centering
  \vspace*{2cm}

  {\color{eyantraBlue}\rule{\linewidth}{3pt}}
  \vspace{0.5cm}

  {\Huge\bfseries\color{eyantraBlue} EEP Auto-Grading Platform}\\[0.6em]
  {\LARGE\color{eyantraAccent} Student Tutorial Guide}\\[0.4em]
  {\large\color{darkgray} E-Yantra Summer Internship Program (EYSIP)}

  \vspace{0.5cm}
  {\color{eyantraBlue}\rule{\linewidth}{3pt}}

  \vspace{1.5cm}

  \begin{tcolorbox}[colback=lightgray, colframe=eyantraAccent, width=0.8\linewidth]
    \centering
    \large This guide walks you through every step of using the\\
    \textbf{EEP Auto-Grading Platform} — from registration to\\
    viewing your final graded results.
  \end{tcolorbox}

  \vfill

  \begin{tabular}{rl}
    \textbf{Document Version:} & 1.0 \\
    \textbf{Platform Version:} & v2.0 \\
    \textbf{Audience:}         & Students (EEP-2 Participants) \\
    \textbf{Organisation:}     & E-Yantra, IIT Bombay \\
  \end{tabular}

  \vspace{1cm}
  {\color{darkgray}\small Last updated: \today}
\end{titlepage}

% ── Table of Contents ─────────────────────────────────────────────────
\tableofcontents
\newpage

% ─────────────────────────────────────────────────────────────────────
\section{Overview}
% ─────────────────────────────────────────────────────────────────────

The \textbf{EEP Auto-Grading Platform} (version 2.0) is a fully online, server-side project evaluation system built for the E-Yantra EEP Software Exploration Program. As a student you can:

\begin{itemize}[leftmargin=*]
  \item Browse and read published weekly assignments
  \item Submit your work via a \textbf{GitHub repository URL} or a \textbf{ZIP file upload}
  \item Watch real-time grading status via live updates
  \item View detailed results including scores, pass/fail status, and execution logs
  \item Attempt quizzes associated with each week's topic
  \item Track your overall progress on the student dashboard
\end{itemize}

\begin{infobox}[How grading works]
  Your submission is executed inside an \emph{isolated Docker sandbox} on our cloud server. The sandbox has \textbf{no internet access} and is destroyed after grading. Results are available within seconds to a few minutes depending on the task complexity.
\end{infobox}

% ─────────────────────────────────────────────────────────────────────
\section{Getting Started}
% ─────────────────────────────────────────────────────────────────────

\subsection{Logging In}

\setcounter{stepnum}{0}

\step{Open the Platform}
Navigate to the platform URL provided by your mentor or program coordinator in your web browser. The login page will be displayed.

\step{Enter Your Credentials}
\begin{itemize}
  \item \textbf{Roll Number / Email:} Your institute roll number (e.g., \texttt{20XXXXX001}) or registered email address.
  \item \textbf{Password:} For newly imported students, the \emph{default password is your email address}. You will be prompted to change it immediately after first login.
\end{itemize}

\step{Change Your Password (First Login)}
\begin{itemize}
  \item On first login, the system enforces a mandatory password change.
  \item Choose a strong password (minimum 8 characters, mix of letters and numbers).
  \item Click \textbf{Update Password} to save.
\end{itemize}

\begin{dangerbox}[Security]
  Do not share your credentials. The grading system logs all submission activity linked to your account.
\end{dangerbox}

\subsection{Dashboard Overview}

After login, you are taken to your \textbf{Student Dashboard}. Here you can see:

\begin{center}
\begin{tabularx}{0.95\linewidth}{>{\bfseries}lX}
\toprule
Widget & Description \\
\midrule
Assignments & Count of published assignments available to you \\
Submissions & Your total submission count across all assignments \\
Passed      & Number of assignments you have passed \\
Avg.\ Score & Your running average score across completed submissions \\
Recent Activity & Latest submissions with status badges \\
\bottomrule
\end{tabularx}
\end{center}

\begin{infobox}
  If you see \textbf{"Classroom approval required"} on the dashboard, your enrollment is still pending. Contact your mentor to approve your classroom enrollment.
\end{infobox}

% ─────────────────────────────────────────────────────────────────────
\section{Browsing Assignments}
% ─────────────────────────────────────────────────────────────────────

\setcounter{stepnum}{0}

\step{Navigate to Assignments}
Click \textbf{Assignments} in the left-hand sidebar. You will see a list of all published weekly assignments.

\step{Open an Assignment}
Click on any assignment card to open its detail page. The detail page contains:

\begin{itemize}
  \item \textbf{Title \& Description} -- What the task is about
  \item \textbf{Category} -- e.g., \texttt{Git}, \texttt{Bash}, \texttt{Python}, \texttt{Docker}
  \item \textbf{Max Score} -- Points available for full completion
  \item \textbf{Deadline} -- Due date/time (UTC)
  \item \textbf{Late Penalty} -- Score deduction percentage applied past the deadline
  \item \textbf{Expected File Structure} -- Required layout for your ZIP submission
  \item \textbf{Submission Instructions} -- Specific instructions to follow
  \item \textbf{Resource Links} -- Learning materials provided by your mentor
\end{itemize}

\begin{warnbox}[Read before submitting]
  Always read the \emph{Expected File Structure} carefully. Submissions that do not match the required layout will fail with a \texttt{VALIDATION\_ERROR}.
\end{warnbox}

% ─────────────────────────────────────────────────────────────────────
\section{Submitting Assignments}
% ─────────────────────────────────────────────────────────────────────

The platform supports two submission methods: \textbf{GitHub URL} and \textbf{ZIP Upload}.

\subsection{Method A: GitHub Repository URL}

\setcounter{stepnum}{0}

\step{Prepare your repository}
Ensure your repository is \textbf{public} on GitHub. Push all your final code before submitting.

\step{Copy the repository URL}
Go to your GitHub repository and copy the HTTPS URL, e.g.:
\begin{lstlisting}
https://github.com/your-username/week1-assignment
\end{lstlisting}

\step{Submit on the platform}
\begin{enumerate}
  \item On the assignment detail page, select \textbf{GitHub URL} as the source type.
  \item Paste the repository URL in the input field.
  \item Click \textbf{Submit}.
\end{enumerate}

\subsection{Method B: ZIP File Upload}

\setcounter{stepnum}{0}

\step{Package your work as a ZIP}
\begin{itemize}
  \item Compress your project folder into a single \texttt{.zip} file.
  \item Ensure the ZIP follows the \emph{Expected File Structure} shown on the assignment page.
  \item Maximum file size: \textbf{50 KB} (50,000 bytes). Keep your ZIP lean — exclude binaries and build artifacts.
\end{itemize}

\step{Upload and submit}
\begin{enumerate}
  \item Select \textbf{ZIP Upload} as the source type.
  \item Click \textbf{Choose File} and select your \texttt{.zip} archive.
  \item Click \textbf{Submit}.
\end{enumerate}

\begin{dangerbox}[ZIP size limit]
  The maximum allowed ZIP size is \textbf{50,000 bytes ($\approx$50 KB)}. Submissions exceeding this will be rejected immediately with a \texttt{FILE\_TOO\_LARGE} error. Do NOT include \texttt{node\_modules}, \texttt{.git}, or compiled binaries in your ZIP.
\end{dangerbox}

\subsection{Submission Rate Limiting}

To prevent accidental double-submissions, the platform enforces a \textbf{5-second cooldown} between submissions for the same assignment. If you hit this limit, wait 5 seconds and try again.

\subsection{Multiple Attempts}

You may re-submit as many times as you wish before the deadline. Each submission is recorded with an \textbf{attempt number}. Your \emph{latest} submission is typically used for final scoring, unless your mentor specifies otherwise.

% ─────────────────────────────────────────────────────────────────────
\section{Tracking Submission Status}
% ─────────────────────────────────────────────────────────────────────

After submitting, the platform updates you in real-time. The submission goes through the following states:

\begin{center}
\begin{tabularx}{0.9\linewidth}{>{\ttfamily\bfseries}l>{\color{darkgray}}X}
\toprule
Status & Meaning \\
\midrule
PENDING          & Submission received, waiting to be queued \\
QUEUED           & Added to the grading job queue \\
RUNNING          & Grader is executing your code in the sandbox \\
COMPLETED        & Grading finished successfully \\
FAILED           & Grader ran but your code did not pass all tests \\
VALIDATION\_ERROR & Submission rejected before grading (wrong format, etc.) \\
TIMEOUT          & Your code exceeded the allowed execution time \\
CANCELLED        & You manually cancelled the submission \\
\bottomrule
\end{tabularx}
\end{center}

\begin{infobox}[Real-time Updates]
  The platform uses \textbf{Server-Sent Events (SSE)} to push live status updates to your browser. You do not need to refresh the page.
\end{infobox}

% ─────────────────────────────────────────────────────────────────────
\section{Viewing Results}
% ─────────────────────────────────────────────────────────────────────

\setcounter{stepnum}{0}

\step{Go to Results}
Click \textbf{Results} in the sidebar, or click on a completed submission card on the dashboard.

\step{Read the Result Detail}
The result page shows:

\begin{itemize}
  \item \textbf{Score} -- Points earned (e.g., \texttt{80 / 100})
  \item \textbf{Passed} -- Green tick for pass, red cross for fail
  \item \textbf{Attempt Number} -- Which attempt this result belongs to
  \item \textbf{Execution Logs (stdout)} -- Output your code produced
  \item \textbf{Error Logs (stderr)} -- Any errors or warnings from the grader
  \item \textbf{AI Feedback} -- Intelligent suggestions to improve your submission (where available)
\end{itemize}

\begin{successbox}[Passed!]
  If your status is \texttt{COMPLETED} and passed is \texttt{true}, your assignment is graded. Celebrate! \faSmile
\end{successbox}

\step{Interpret a Failed Submission}
\begin{enumerate}
  \item Read the \textbf{stderr} output carefully — it often shows which test case failed and why.
  \item Fix the issue in your code.
  \item Re-submit before the deadline.
\end{enumerate}

% ─────────────────────────────────────────────────────────────────────
\section{Quizzes}
% ─────────────────────────────────────────────────────────────────────

Each week may include a quiz to test your conceptual understanding.

\setcounter{stepnum}{0}

\step{Navigate to Quizzes}
Click \textbf{Quizzes} or find the quiz linked from the assignment page.

\step{Start the Quiz}
\begin{itemize}
  \item Quizzes are timed. A countdown timer is shown at the top.
  \item Questions may be multiple-choice (single or multiple answers) or short text.
  \item Once you start, the timer cannot be paused.
\end{itemize}

\step{Submit your Quiz}
\begin{itemize}
  \item Answer all questions and click \textbf{Submit Quiz}.
  \item Results are shown immediately after submission.
\end{itemize}

\begin{warnbox}[Quiz rules]
  \begin{itemize}
    \item Do not refresh the page during a quiz -- your progress may be lost.
    \item Quizzes can only be attempted once unless your mentor resets them.
    \item Each quiz has a fixed time window; once the deadline passes, you cannot start.
  \end{itemize}
\end{warnbox}

% ─────────────────────────────────────────────────────────────────────
\section{Announcements}
% ─────────────────────────────────────────────────────────────────────

Mentors post important notices in the \textbf{Announcements} section. You will receive in-platform notifications when a new announcement is posted.

\begin{itemize}
  \item Click \textbf{Announcements} in the sidebar to view all notices.
  \item Read each announcement before making submissions -- mentors sometimes clarify requirements or extend deadlines here.
\end{itemize}

% ─────────────────────────────────────────────────────────────────────
\section{Profile \& Settings}
% ─────────────────────────────────────────────────────────────────────

\begin{itemize}
  \item Click your name in the top-right corner and select \textbf{Profile}.
  \item You can update your \textbf{full name} and \textbf{email address}.
  \item To change your password, use the \textbf{Change Password} form on the profile page.
\end{itemize}

% ─────────────────────────────────────────────────────────────────────
\section{Common Issues \& Solutions}
% ─────────────────────────────────────────────────────────────────────

\begin{center}
\begin{tabularx}{\linewidth}{>{\bfseries}lX}
\toprule
Problem & Solution \\
\midrule
"Classroom approval required" & Contact your mentor to approve your enrollment. \\
ZIP rejected (too large)       & Remove \texttt{node\_modules}, \texttt{.git}, compiled files. \\
VALIDATION\_ERROR               & Check the expected file structure on the assignment page. \\
TIMEOUT                        & Optimise your code. Avoid infinite loops. \\
GitHub URL not accepted        & Ensure the repository is \textbf{public}. \\
Submission stuck in QUEUED     & The grading worker may be busy. Wait a few minutes. \\
Login fails                    & Use your roll number (not email) as the username. \\
\bottomrule
\end{tabularx}
\end{center}

% ─────────────────────────────────────────────────────────────────────
\section{Quick Reference}
% ─────────────────────────────────────────────────────────────────────

\begin{tcolorbox}[colback=eyantraBlue!5, colframe=eyantraBlue, title={\faBookmark\enspace Student Quick Reference}]
\begin{tabular}{ll}
\textbf{Login credential:}   & Roll number or email \\
\textbf{Default password:}   & Your email address (change immediately) \\
\textbf{Max ZIP size:}        & 50,000 bytes ($\approx$50 KB) \\
\textbf{Rate limit:}          & 5 seconds between submissions \\
\textbf{Sandbox network:}    & None (no internet inside grading container) \\
\textbf{API Docs (debug):}   & \texttt{/docs} (if enabled by admin) \\
\end{tabular}
\end{tcolorbox}

% ─────────────────────────────────────────────────────────────────────
\end{document}
